\documentclass[11pt]{article}
\usepackage[margin=0.85in]{geometry}
\usepackage[T1]{fontenc}
\usepackage{amsmath,amssymb,booktabs,graphicx,xurl,hyperref}
\graphicspath{{outputs/task1/}{outputs/task2/}{outputs/task3/debug/}{outputs/task3/lego_train/}{outputs/task3/lego_train/validation/}{outputs/task3/lego_train/mesh/}}

\title{CS 498 HW4 Report\\Tiny NeRF for Novel View Synthesis}
\author{Your Name (NetID)}
\date{Fall 2026}

\begin{document}
\maketitle

\section{Rays and positional encoding}
Describe the camera-coordinate convention and the world-space ray equation.
Explain why image rows reverse camera $y$ and why ray directions remain
unnormalized until rendering. Include
\path{task1_visual_checkpoint.png}. State the encoded dimensions for the
provided full-run position and direction settings.

\section{Tiny NeRF and volume rendering}
Describe the provided network architecture and your forward pass, including
the skip connection, density head, and view-dependent color branch. Give the
sampling and alpha-compositing equations you implemented. Include
\path{task2_visual_checkpoint.png} and briefly interpret its deterministic
depth, perturbed depth, and compositing panels.

\section{Training and held-out reconstruction}
Report whether the smoke run completed. Include
\path{train_psnr_curve.png} and the two final milestone renders
\path{val_010000_view_00.png} and \path{val_010000_view_01.png}. State the
final logged training PSNR.

Report the mean validation PSNR and SSIM from
\path{validation_metrics.csv}. Include at least one
\path{view_00_comparison.png} from the validation directory and
\path{turntable_contact_sheet.png}. Describe one visible limitation or
view-dependent failure.

\section{Density grid and textured mesh}
State the density-grid bounds, resolution, threshold, mesh vertex/face counts,
and texture-atlas size. Include the two
\path{view_00_comparison.png}/\path{view_01_comparison.png} figures from the
mesh directory. Report the mean direct-NeRF and textured-mesh PSNR/SSIM values
from \path{mesh/comparison.csv}. Explain two visible differences between
direct volume rendering and textured-mesh rendering.

\section{Limitations and reproducibility}
Summarize the commands and hardware used for the full run. Discuss at least
two limitations of this compact NeRF or its mesh extraction.

\section*{Optional ball-insertion bonus}
If attempted, state the supplied sphere parameters and explain the
transmittance-aware compositing rule you implemented. Include
\path{ball_insertion_contact_sheet.png} and
\path{ball_visibility.png}. Identify two views with different occlusion
relationships and explain why expected-depth z-buffering is insufficient.
Remove this section if you did not attempt the bonus.

\section*{AI-use declaration}
Check one statement.

\begin{itemize}
  \item[$\square$] I did not use generative AI for this homework.
  \item[$\square$] I used generative AI only in ways explicitly allowed by
  the course policy. Tool/model, purpose, and affected files:
  \underline{\hspace{1.6in}}
\end{itemize}

I confirm that the submitted code, outputs, numerical results, and written
analysis comply with the course collaboration and AI-use policy.

\end{document}
